APR 04 - TEORIA OTIMIZACAO/DEEP
Rede profunda tem varias camadas; aprende representacoes hierarquicas.
Problemas: custo computacional, muitos parametros, overfitting e gradientes que desaparecem/explodem.
Backprop calcula gradientes; optimizer define como atualizar parametros.
GD batch: gradiente dataset inteiro; SGD: por amostra; mini-batch: blocos intermedios.
Learning rate muito alto pode divergir; muito baixo e lento.
Momentum suaviza oscilacoes usando historico de gradientes.
RMSProp ajusta passo por parametro a partir dos quadrados recentes.
Adam combina momentum (primeiro momento) e escala adaptativa (segundo momento), com correcao de bias.
ReLU: reduz saturacao para z>0 e e simples; leaky ReLU ajuda contra dying ReLU.
Sigmoid e tanh podem saturar; derivadas pequenas nas caudas.
Vanishing gradient: primeiras camadas recebem gradientes quase zero; exploding gradient: magnitude cresce.
Regularizacao L1 favorece esparsidade; L2 penaliza pesos grandes; dropout reduz coadaptacao.
Data augmentation cria variacoes de treino relevantes; nao alterar semantica dos labels.
Early stopping usa validacao, nunca conjunto de teste.
Batch normalization normaliza ativacoes dentro de lotes e tem parametros treinaveis.
Deep autoencoder: encoder comprime em representacao latente, decoder reconstrui input; usado para reducao, features ou anomalias.
Transfer learning reutiliza pesos de modelo pre-treinado e ajusta camada(s).
Parametros sao W,b; hiperparametros sao eta, arquitetura, batch size, epocas, lambda.
Erro treino baixo e val alto => overfit. Ambos altos => underfit.
Para classificar: sigmoid+BCE ou softmax+cross entropy, conforme classes.